\documentclass[aspectratio=169, xcolor=dvipsnames]{beamer}

% --- Theme Setup ---
\usetheme{Madrid}
\usecolortheme{default}

% --- Packages ---
\usepackage{graphicx}
\usepackage{xcolor}
\usepackage{booktabs}
\usepackage{hyperref}
\usepackage{adjustbox}
\usepackage{amssymb}
\usepackage{longtable}
\usepackage{caption}

% --- Title Information ---
\title{Module 03}
\subtitle{Why DBMS?/2}
\author{Milav Dabgar}
\institute{www.milav.in}
\date{\today}

\begin{document}

% Title Slide
\begin{frame}
    \titlepage
\end{frame}

\begin{frame}{Module Recap}
    \begin{itemize}
        \item Evolution of Data and Records Management
        \item History of DBMS
    \end{itemize}
\end{frame}

\begin{frame}{Module Objectives}
    \begin{itemize}
        \item Comparison of File based data management and DBMS
    \end{itemize}
\end{frame}

\begin{frame}{Module Outline}
    \begin{itemize}
        \item File handling by Python viz-a-viz DBMS - Bank Transaction example
        \item Parameterized Comparison
    \end{itemize}
\end{frame}

\section{File Systems vs Databases}

\begin{frame}
  \vfill
  \centering
  \begin{beamercolorbox}[sep=8pt,center,shadow=true,rounded=true]{title}
    \usebeamerfont{title}File Systems vs Databases\par%
  \end{beamercolorbox}
  \vfill
\end{frame}

\begin{frame}{Case study: A bank transaction}
    \textbf{Banking Transaction System} \\
    Consider a simple banking system where a person can open a new account, transfer fund to an existing account and check the history of all her transactions till date.
    
    \vspace{1em}
    The application performs the following checks:
    \begin{itemize}
        \item If the account balance is not enough, it will not allow the fund transfer
        \item If the account numbers are not correct, it will flash a message and terminate the transaction.
        \item If a transaction is successful, it prints a confirmation message.
    \end{itemize}
\end{frame}

\begin{frame}{Case study: A bank transaction (2)}
    We will use this banking transaction system to compare various features of a file-based (spreadsheet/.csv files) implementation viz-a-viz a DBMS-based implementation
    
    \begin{itemize}
        \item Account details are stored in
        \begin{itemize}
            \item \texttt{Accounts.csv} for file-based implementation
            \item \texttt{Accounts} table for DBMS implementation
        \end{itemize}
        \item The transaction details are stored in
        \begin{itemize}
            \item \texttt{Ledger.csv} file for file-based implementation
            \item \texttt{Ledger} table for DBMS implementation
        \end{itemize}
    \end{itemize}
    
    \vspace{1em}
    In the following slides we discuss a fund transfer transaction.\\
    \textbf{Source:} \url{https://github.com/bhaskariitm/transition-from-files-to-db/tree/main}
\end{frame}

\section{Python viz-a-viz SQL}

\begin{frame}
  \vfill
  \centering
  \begin{beamercolorbox}[sep=8pt,center,shadow=true,rounded=true]{title}
    \usebeamerfont{title}Python viz-a-viz SQL\par%
  \end{beamercolorbox}
  \vfill
\end{frame}

\begin{frame}[fragile]{Bank Transaction: Python viz-a-viz SQL}
    \begin{columns}[T]
        \begin{column}{0.5\textwidth}
            \textbf{Python}
            \begin{verbatim}
def begin_Transaction(creditAcc, debitAcc, amount):
    temp = [] 
    success = 0
    # Open file handles to retrieve and store transaction data
    f_obj_Account1 = open('Accounts.csv', 'r')
    f_reader1 = csv.DictReader(f_obj_Account1)
    
    f_obj_Account2 = open('Accounts.csv', 'r')
    f_reader2 = csv.DictReader(f_obj_Account2)
    
    f_obj_Ledger = open('Ledger.csv', 'a+')
    f_writer = csv.DictWriter(f_obj_Ledger, fieldnames=col_name_Ledger)
            \end{verbatim}
        \end{column}
        \begin{column}{0.5\textwidth}
            \textbf{SQL} \\
            \vspace{1em}
            \textit{// Handled implicitly by the DBMS}
        \end{column}
    \end{columns}
\end{frame}

\begin{frame}[fragile]{Bank Transaction: Python viz-a-viz SQL (2)}
    \begin{columns}[T]
        \begin{column}{0.5\textwidth}
            \textbf{Python}
            \begin{verbatim}
try:
    for sRec in f_reader1: 
        #CONDITION CHECK FOR ENOUGH BALANCE
        if sRec["AcctNo"] == debitAcc and int(sRec["Balance"]) > int(amt):
            ...
            \end{verbatim}
        \end{column}
        \begin{column}{0.5\textwidth}
            \textbf{SQL}
            \begin{verbatim}
do $$
begin
    amt = 5000;
    sendVal = '1800090';
    recVal = '1800100';
    
    select balance from accounts 
        into sbalance 
        where account_no = sendVal;
        
    if sbalance < amt then
        ...
$$
            \end{verbatim}
        \end{column}
    \end{columns}
\end{frame}

\begin{frame}[fragile]{Bank Transaction: Python viz-a-viz SQL (3)}
    \begin{columns}[T]
        \begin{column}{0.5\textwidth}
            \textbf{Python}
            \begin{verbatim}
# in continuation...
        for rRec in f_reader2:
            if rRec["AcctNo"] == creditAcc:
                sRec["Balance"] = str(int(sRec["Balance"]) - int(amt)) # DEBIT
                temp.append(sRec)
                ...
            \end{verbatim}
        \end{column}
        \begin{column}{0.5\textwidth}
            \textbf{SQL}
            \begin{verbatim}
    # in continuation...
    if sbalance < amt then
        raise notice "Insufficient balance";
    else
        update accounts 
            set balance = balance - amt 
            where account_no = sendVal;
        ...
$$
            \end{verbatim}
        \end{column}
    \end{columns}
\end{frame}

\begin{frame}[fragile]{Bank Transaction: Python viz-a-viz SQL (4)}
    \begin{columns}[T]
        \begin{column}{0.5\textwidth}
            \textbf{Python}
            \begin{verbatim}
# in continuation...
                # Critical point
                f_writer.writerow({
                    "Acct1":sRec["AcctNo"],
                    "Acct2":rRec["AcctNo"],
                    "Amount":amt, "D/C":"D"})
                
                rRec["Balance"] = str(int(rRec["Balance"]) + int(amt)) # CREDIT
                temp.append(rRec)
                ...
            \end{verbatim}
        \end{column}
        \begin{column}{0.5\textwidth}
            \textbf{SQL}
            \begin{verbatim}
# in continuation...
        insert into ledger(sendAc, recAc, amnt, ttype) 
            values(sendVal, recVal, amt, 'D');
            
        update accounts 
            set balance = balance + amt 
            where account_no = recVal;
        ...
$$
            \end{verbatim}
        \end{column}
    \end{columns}
\end{frame}

\begin{frame}[fragile]{Bank Transaction: Python viz-a-viz SQL (5)}
    \begin{columns}[T]
        \begin{column}{0.5\textwidth}
            \textbf{Python}
            \begin{verbatim}
# in continuation...
                f_writer.writerow({
                    "Acct1":rRec["AcctNo"],
                    "Acct2":sRec["AcctNo"],
                    "Amount":amt,"D/C":"C"})
                success = success + 1
                break
                
        f_obj_Account1.seek(0)
        next(f_obj_Account1)
        for record in f_reader1:
            if record["AcctNo"] != temp[0]["AcctNo"] and record["AcctNo"] != temp[1]["AcctNo"]:
except: 
    print("Wrong input entered !!!")
            \end{verbatim}
        \end{column}
        \begin{column}{0.5\textwidth}
            \textbf{SQL}
            \begin{verbatim}
# in continuation...
        insert into ledger(sendAc, recAc, amnt, ttype) 
            values(recVal, sendVal, amt, 'C');
            
        commit;
        raise notice "Successful";
    end if;
end;
$$
            \end{verbatim}
        \end{column}
    \end{columns}
\end{frame}

\begin{frame}[fragile]{Bank Transaction: Python viz-a-viz SQL (6)}
    \begin{columns}[T]
        \begin{column}{0.5\textwidth}
            \textbf{Python}
            \begin{verbatim}
#Writing back to the file 
f_obj_Account1.close()
f_obj_Account2.close()
f_obj_Ledger.close()

if success == 1:
    f_obj_Account = open('Accounts.csv', 'w+', newline='')
    f_writer = csv.DictWriter(f_obj_Account, fieldnames=col_name_Account)
    f_writer.writeheader()
    for data in temp:
        f_writer.writerow(data)
    f_obj_Account.close()
    print("Transaction is successful !!")
else:
    print('Transaction failed: Confirm Account details')
            \end{verbatim}
        \end{column}
        \begin{column}{0.5\textwidth}
            \textbf{SQL} \\
            \vspace{1em}
            \textit{// Handled implicitly by the DBMS}
        \end{column}
    \end{columns}
\end{frame}

\section{Parameterized Comparison}

\begin{frame}
  \vfill
  \centering
  \begin{beamercolorbox}[sep=8pt,center,shadow=true,rounded=true]{title}
    \usebeamerfont{title}Parameterized Comparison\par%
  \end{beamercolorbox}
  \vfill
\end{frame}

\begin{frame}{Comparison}
    \begin{center}
        \begin{adjustbox}{max width=\textwidth}
        \begin{tabular}{|l|p{6cm}|p{6cm}|}
        \hline
        \textbf{Parameter} & \textbf{File Handling via Python} & \textbf{DBMS} \\ \hline
        Scalability (data amount) & Very difficult to handle insert, update and querying of records & In-built features to provide high scalability for a large number of records \\ \hline
        Scalability (structure change) & Extremely difficult to change the structure of records as in the case of adding or removing attributes & Adding or removing attributes can be done seamlessly using simple SQL queries \\ \hline
        Time of execution & In seconds & In milliseconds \\ \hline
        Persistence & Data processed using temporary data structures have to be manually updated to the file & Data persistence is ensured via automatic, system induced mechanisms \\ \hline
        Robustness & Ensuring robustness of data has to be done manually & Backup, recovery and restore need minimum manual intervention \\ \hline
        Security & Difficult to implement in Python (Security at OS level) & User-specific access at database level \\ \hline
        Programmer's productivity & Most file access operations involve extensive coding to ensure persistence, robustness and security of data & Standard and simple built-in queries reduce the effort involved in coding thereby increasing a programmer's throughput \\ \hline
        Arithmetic operations & Easy to do arithmetic computations & Limited set of arithmetic operations are available \\ \hline
        Costs & Low costs for hardware, software and human resources & High costs for hardware, software and human resources \\ \hline
        \end{tabular}
        \end{adjustbox}
    \end{center}
\end{frame}

\begin{frame}{Scalability}
    \textbf{File handling via Python}
    \begin{itemize}
        \item \textbf{Number of records:} As the \# of records increases, the efficiency of flat files reduces:
        \begin{itemize}
            \item the time spent in searching for the right records
            \item the limitations of the OS in handling huge files
        \end{itemize}
        \item \textbf{Structural Change:} To add an attribute, initializing the new attribute of each record with a default value has to be done by program. It is very difficult to detect and maintain relationships between entities if and when an attribute has to be removed.
    \end{itemize}
    
    \vspace{1em}
    \textbf{DBMS}
    \begin{itemize}
        \item \textbf{Number of records:} Databases are built to efficiently scale up when the \# of records increase drastically.
        \begin{itemize}
            \item In-built mechanisms, like indexing, for quick access of right data.
        \end{itemize}
        \item \textbf{Structural Change:} During adding an attribute, a default value can be defined that holds for all existing records - the new attribute gets initialized with the default value. During deletion, constraints are used either not to allow the removal or ensure its safe removal
    \end{itemize}
\end{frame}

\begin{frame}{Time and Efficiency}
    \textbf{File handling via Python}
    \begin{itemize}
        \item The effort needed to implement a file handler is quite less in Python
        \item In order to process a 1GB file, a program in Python would typically take few seconds.
    \end{itemize}
    
    \vspace{1em}
    \textbf{DBMS}
    \begin{itemize}
        \item The effort to install and configure a DB in a DB server is expensive \& time consuming
        \item In order to process a 1GB file, an SQL query would typically take few milliseconds.
        \item If the number of records is very small, the overhead in installing and configuring a database will be much more than the time advantage obtained from executing the queries.
        \item However, if the number of records is really large, then the time required in the initialization process of a database will be negligible as compared to the time saved in using SQL queries.
    \end{itemize}
\end{frame}

\begin{frame}{Persistence, Robustness, Security}
    \textbf{File handling via Python}
    \begin{itemize}
        \item \textbf{Persistence:} Data processed using in-memory data structures stay in the memory during processing. After updates, these are manually updated to the file on disk
        \item \textbf{Robustness:} Ensuring consistency, reliability and sanity is manual via multiple checks. On a system crash, a transaction may cause inconsistency or loss of data.
        \item \textbf{Security:} Extremely difficult to implement granular security in file systems. Authentication is at the OS level.
    \end{itemize}
    
    \vspace{1em}
    \textbf{Database Management Systems (DBMS)}
    \begin{itemize}
        \item \textbf{Persistence:} Data persistence is ensured via automatic, system mechanisms. The programmer does not have to worry about the data getting lost due to manual errors
        \item \textbf{Robustness:} Backup, recovery \& restore need minimum manual intervention. The backup and recovery plan can be devised for automatic recovery on a crash
        \item \textbf{Security:} DBMS provides user-specific access at the database level with restriction for to view only access
    \end{itemize}
\end{frame}

\begin{frame}{Programmer's Productivity}
    \textbf{File handling via Python}
    \begin{itemize}
        \item \textbf{Building the file handler:} Since the constraints within and across entities have to be enforced manually, the effort involved in building a file handling application is huge
        \item \textbf{Maintenance:} To maintain the consistency of data, one must regularly check for sanity of data and the relationships between entities during inserts, updates and deletes
        \item \textbf{Handling huge data:} As the data grows beyond the capacity of the file handler, more efforts are needed
    \end{itemize}
    
    \vspace{1em}
    \textbf{DBMS}
    \begin{itemize}
        \item \textbf{Configuring the database:} The installation and configuration of a database is specialized job of a DBA. A programmer, on the other hand, is saved the trouble
        \item \textbf{Maintenance:} DBMS has in-built mechanisms to ensure consistency and sanity of data being inserted, updated or deleted. The programmer does not need to do such checks
        \item \textbf{Handling huge data:} DBMS can handle even terabytes of data - Programmer does not have to worry
    \end{itemize}
\end{frame}

\begin{frame}{Arithmetic Operations}
    \textbf{File handling via Python}
    \begin{itemize}
        \item \textbf{Extensive support for arithmetic and logical operations:} Extensive arithmetic and logical operations can be performed on data using Python. These include complex numerical calculations and recursive computations.
    \end{itemize}
    
    \vspace{1em}
    \textbf{DBMS}
    \begin{itemize}
        \item \textbf{Limited support for arithmetic and logical operations:} SQL provides limited arithmetic and logical operations. Any other complex computation has to be done outside the SQL.
    \end{itemize}
\end{frame}

\begin{frame}{Costs and Complexity}
    \textbf{File handling via Python}
    \begin{itemize}
        \item File systems are cheaper to install and use. No specialized hardware, software or personnel are required to maintain filesystems.
    \end{itemize}
    
    \vspace{1em}
    \textbf{DBMS}
    \begin{itemize}
        \item Large databases are served by dedicated database servers need large storage and processing power
        \item DBMSs are expensive software that have to be installed and regularly updated
        \item Databases are inherently complex and need specialized people to work on it like \textbf{DBA}
        \item The above factors lead to \textbf{huge costs in implementing and maintaining database management systems}
    \end{itemize}
\end{frame}

\begin{frame}{Module Summary}
    \begin{itemize}
        \item Elucidated the difference between File handling by Python viz-a-viz DBMS through an Bank Transaction example
        \item Parameterized Comparison
    \end{itemize}
    
    \vspace{2em}
    \begin{center}
    \small \textit{Slides used in this presentation are borrowed from \url{http://db-book.com/} with kind permission of the authors.} \\
    \small \textit{Edited and new slides are marked with 'PPD'.}
    \end{center}
\end{frame}

\end{document}
